% =====================================================================
%  CARNET D'ENTRAÎNEMENT — Fiche 12 (Chimie organique)
%  THÈME 1 — Hydrocarbures, formules brutes et squelette carboné
%  NIVEAU MOYEN — ÉNONCÉS
% =====================================================================
\documentclass[11pt,a4paper]{article}
\usepackage[utf8]{inputenc}
\usepackage[T1]{fontenc}
\usepackage{lmodern}
\usepackage[french]{babel}
\usepackage{amsmath,amssymb,mathtools}
\providecommand{\degree}{^\circ}
\usepackage{icomma}
\usepackage{siunitx}
\sisetup{output-decimal-marker={,},per-mode=symbol}
\usepackage[version=4]{mhchem}
\usepackage{chemfig}
\setchemfig{atom sep=2em,bond length=0.9em,cram width=2.5pt}
\usepackage{xcolor}
\usepackage{array}
\usepackage{booktabs}
\usepackage{tabularx}
\usepackage[most]{tcolorbox}
\usepackage{enumitem}
\usepackage{tikz}
\usetikzlibrary{arrows.meta,positioning,calc}
\usepackage{fancyhdr}
\usepackage{titlesec}
\usepackage[hidelinks]{hyperref}
\usepackage[a4paper,margin=2.1cm,headheight=26pt,headsep=9pt]{geometry}

\definecolor{primary}{RGB}{150,98,162}
\definecolor{primarydark}{RGB}{92,52,110}
\definecolor{excol}{RGB}{0,135,90}
\definecolor{attncol}{RGB}{200,40,55}
\definecolor{methcol}{RGB}{90,90,100}

\newcommand{\runtitle}{Thème 1 — Moyen — Énoncés}
\pagestyle{fancy}\fancyhf{}
\fancyhead[L]{\small\textcolor{primarydark}{\textbf{Fiche 12} — Chimie organique}}
\fancyhead[R]{\small\textcolor{primarydark}{\runtitle}}
\fancyfoot[C]{\small\thepage}
\renewcommand{\headrulewidth}{0.8pt}
\renewcommand{\headrule}{\hbox to\headwidth{\color{primary}\leaders\hrule height \headrulewidth\hfill}}
\setlength{\parindent}{0pt}\setlength{\parskip}{3pt}

\tcbset{boxbase/.style={enhanced,breakable,boxrule=0.9pt,arc=2.5pt,
   left=3mm,right=3mm,top=2mm,bottom=2mm,fonttitle=\bfseries,coltitle=white,
   toptitle=0.5mm,bottomtitle=0.5mm}}
\newtcolorbox[auto counter]{exo}[1]{boxbase,colback=primary!4,colframe=primary,
   title={\textsc{Exercice}~\thetcbcounter\ \textmd{--- \itshape #1}}}
\newtcolorbox{consigne}{boxbase,colback=methcol!8,colframe=methcol,sharp corners,
   title={\textsc{Consignes}}}

\newcommand{\banner}[2]{%
\begin{tcolorbox}[enhanced,colback=primary,colframe=primary,arc=5pt,boxrule=0pt,
   left=5mm,right=5mm,top=2.5mm,bottom=2.5mm,coltext=white]
{\large\bfseries #1}\\[1pt]{\small #2}
\end{tcolorbox}\vspace{2pt}}

\begin{document}

\banner{Thème 1 — Hydrocarbures, formules brutes et squelette}{Niveau \textbf{Moyen} — Énoncés \quad$\vert$\quad Fiche 12, Chimie organique}

\begin{consigne}
Exercices \textbf{ouverts}. On combine ici plusieurs mécaniques du tuto : degré d'insaturation, comptage des H par carbone, isomérie de constitution et stéréoisomérie cis/trans. Justifie toujours par le décompte des liaisons.
\end{consigne}

% ---------------------------------------------------------------------
\begin{exo}{du degré d'insaturation aux structures possibles}
Une molécule a pour formule brute $\ce{C5H8}$.
\begin{enumerate}[label=\textbf{\alph*)},leftmargin=2em,itemsep=1pt]
  \item Calcule son indice de déficience en hydrogène $I$.
  \item Que représente concrètement la valeur trouvée (liaisons $\pi$ et/ou cycles) ?
  \item Propose \textbf{trois} structures différentes (familles ou arrangements distincts) compatibles avec cette formule brute.
\end{enumerate}
\end{exo}

% ---------------------------------------------------------------------
\begin{exo}{trier les carbones par nombre d'hydrogènes}
On considère l'hydrocarbure ci-dessous (\textbf{4-méthylpent-1-ène}) :
\begin{center}\chemfig{CH_2=[:30]CH-[:-30]CH_2-[:30]CH(-[:90]CH_3)-[:-30]CH_3}\end{center}
\begin{enumerate}[label=\textbf{\alph*)},leftmargin=2em,itemsep=1pt]
  \item Détermine sa formule brute et vérifie-la par le calcul de $I$.
  \item Indique combien de carbones portent respectivement \textbf{3 H}, \textbf{2 H}, \textbf{1 H} et \textbf{0 H}. (Attention : un carbone de double liaison forme encore d'autres liaisons.)
\end{enumerate}
\end{exo}

% ---------------------------------------------------------------------
\begin{exo}{la stéréoisomérie cis / trans}
On s'intéresse au \textbf{but-2-ène} $\ce{CH3-CH=CH-CH3}$.
\begin{enumerate}[label=\textbf{\alph*)},leftmargin=2em,itemsep=1pt]
  \item Pourquoi le but-2-ène existe-t-il sous deux formes géométriques distinctes, alors que le but-1-ène $\ce{CH2=CH-CH2-CH3}$ n'en a qu'une ?
  \item Fais un schéma des configurations \emph{cis} et \emph{trans} et précise dans chaque cas la disposition des deux groupes $\ce{CH3}$.
  \item Ces deux formes ont-elles la même formule brute ? Sont-elles pour autant la même molécule ? Justifie.
\end{enumerate}
\end{exo}

% ---------------------------------------------------------------------
\begin{exo}{les trois isomères de \ce{C5H12}}
La formule brute $\ce{C5H12}$ correspond à trois isomères de constitution.
\begin{enumerate}[label=\textbf{\alph*)},leftmargin=2em,itemsep=1pt]
  \item Dessine (formule topologique ou semi-développée) et nomme ces \textbf{trois} isomères.
  \item Pour chaque isomère, donne le décompte des carbones \emph{primaires / secondaires / tertiaires / quaternaires}.
\end{enumerate}
\end{exo}

% ---------------------------------------------------------------------
\begin{exo}{reconstruire une molécule à partir de contraintes}
Un \textbf{alcane} a pour formule brute $\ce{C6H14}$. On sait qu'il ne contient \textbf{aucun carbone secondaire}, qu'il possède exactement \textbf{deux carbones tertiaires} (liés l'un à l'autre) et \textbf{quatre carbones primaires}.
\begin{enumerate}[label=\textbf{\alph*)},leftmargin=2em,itemsep=1pt]
  \item Déduis-en la structure de la molécule et dessine-la.
  \item Nomme-la et vérifie que le décompte des carbones et des hydrogènes redonne bien $\ce{C6H14}$.
\end{enumerate}
\end{exo}

\end{document}
